\VerdierPhysicalPageAnchor{0113}%
\begingroup
\small
\setlength{\abovedisplayskip}{4pt plus 2pt minus 2pt}%
\setlength{\belowdisplayskip}{4pt plus 2pt minus 2pt}%

\VerdierSemanticAnchor{ch2-1-1-5-additivity-announcement}
\noindent On montrera (1.2.7) que les foncteurs cohomologiques sont
additifs et que les foncteurs multi-exacts sont multi-additifs.

\VerdierSemanticAnchor{ch2-1-1-5-cohomology-notation}
Soit $H\colon\mathcal D\longrightarrow\mathcal A$ un foncteur
cohomologique. On utilise le plus souvent la notation
$H^0(X)=H(X)$, $H^n(X)=H(X[n])$. En utilisant l'axiome TRII, on voit
qu'un foncteur cohomologique associe à tout triangle distingué

\VerdierSemanticAnchor{ch2-1-1-5-long-exact-setup}
\[
  X\xrightarrow{u}Y\xrightarrow{v}Z\xrightarrow{w}X[1]
\]
une suite exacte longue :

\VerdierSemanticAnchor{ch2-1-1-5-1}
\begin{equation}
  \cdots\longrightarrow H^{-1}(Z)
  \xrightarrow{H^{-1}(-w)}H^0(X)
  \xrightarrow{H^0(u)}H^0(Y)
  \xrightarrow{H^0(v)}H^0(Z)
  \xrightarrow{H^0(w)}\cdots
  \tag{1.1.5.1}
\end{equation}

\VerdierSemanticAnchor{ch2-1-1-6}
\medskip
\noindent\textbf{1.1.6.}\quad
Soit $\mathcal D$ une catégorie triangulée. On emploie le plus souvent
la terminologie suivante : un triangle

\VerdierSemanticAnchor{ch2-1-1-6-antidistinguished-triangle}
\[
  X\xrightarrow{u}Y\xrightarrow{v}Z\xrightarrow{w}X[1]
\]
de $\mathcal D$ est dit \emph{antidistingué} si le triangle
\[
  X\xrightarrow{u}Y\xrightarrow{v}Z\xrightarrow{-w}X[1]
\]
est distingué.

\VerdierSemanticAnchor{ch2-1-1-6-antiexact-functor}
Un $\mathbb Z$-foncteur d'une catégorie triangulée $\mathcal D'$ dans
$\mathcal D$ est dit \emph{antiexact} s'il transforme les triangles
distingués de $\mathcal D'$ en triangles antidistingués de $\mathcal D$.

\VerdierSemanticAnchor{ch2-1-1-6-exact-functor-notation}
Le plus souvent, pour désigner un foncteur exact de $\mathcal D'$ dans
$\mathcal D$, on omet de faire figurer dans la notation les isomorphismes
de commutation aux translations.

\VerdierSemanticAnchor{ch2-1-1-7}
\medskip
\noindent\textbf{1.1.7.}\quad
Soient $\mathcal C$ une $\mathbb Z$-catégorie stricte (chap. I, 1.2.2)
et :

\VerdierSemanticAnchor{ch2-1-1-7-translations}
\[
  X\longmapsto X[n],\qquad n\in\mathbb Z,
\]
les foncteurs de translation de $\mathcal C$. Soit $\mathcal C^\circ$ la
$\mathbb Z$-catégorie opposée (chap. I, 1.3.9).

\VerdierSemanticAnchor{ch2-1-1-7-opposite-morphism}
Pour tout morphisme $u$ de $\mathcal C$, notons $((u))$ le morphisme
correspondant de $\mathcal C^\circ$. Cette notation est abusive car
l'ensemble des objets et l'ensemble des morphismes de $\mathcal C^\circ$
sont égaux respectivement à l'ensemble des objets et l'ensemble des
morphismes de $\mathcal C$. Désignons par :

\VerdierSemanticAnchor{ch2-1-1-7-opposite-translations}
\[
  X\longmapsto X[[n]],\qquad n\in\mathbb Z,
\]
les foncteurs de translations de $\mathcal C^\circ$.

\VerdierSemanticAnchor{ch2-1-1-7-translation-identity}
Rappelons qu'on a, pour tout objet $X$ de $\mathcal C$ et tout entier
$n\in\mathbb Z$ (chap. I, 1.3.9) :
\[
  X[-n]=X[[n]].
\]

\VerdierSemanticAnchor{ch2-1-1-7-opposite-category}
On appelle \emph{catégorie triangulée opposée} à une catégorie triangulée
$\mathcal D$ la $\mathbb Z$-catégorie stricte $\mathcal D^\circ$ opposée
à la catégorie $\mathcal D$ (chap. I, 1.3.9) munie de l'ensemble de
triangles distingués suivant. Un triangle de $\mathcal D^\circ$ :

\VerdierSemanticAnchor{ch2-1-1-7-opposite-triangle}
\[
  X[[-1]]\xrightarrow{((w))}Z\xrightarrow{((v))}Y
  \xrightarrow{((u))}X
\]
est distingué si et seulement si le triangle correspondant de la catégorie
$\mathcal D$ :

\VerdierSemanticAnchor{ch2-1-1-7-corresponding-antitriangle}
\[
  X\xrightarrow{u}Y\xrightarrow{v}Z\xrightarrow{w}X[1]
\]
est \emph{antidistingué}.

\VerdierSemanticAnchor{ch2-1-1-7-triangulated-axioms}
On démontre aisément que la famille de triangles de $\mathcal D^\circ$
ainsi définie vérifie les axiomes des catégories triangulées.
\endgroup

